\chapter{Social Reform in Indian Society}

\section{Objectives and Social Composition of the Reformers}

Most social reforms in India are traced to \textbf{1830 onwards}, and what was reformed largely concerned \textbf{caste practices, the status of women, and untouchability}. (Note: UPSC has not asked questions from 'protest and movement during the colonial period' --- Section A --- but has asked from 'social movements' in Section C, which is treated there.)

\begin{KeyConcept}
\begin{itemize}
\item The \textbf{three broad objectives of social reform}, common across all reformers:
\item 1. \textbf{Emancipation of women}.
\item 2. \textbf{Abolition of untouchability}.
\item 3. \textbf{Reform of the caste structure}.
\item Whether it is the Ramakrishna Mission, Swami Vivekananda, Periyar, Jyotiba Phule, Savitribai Phule, Mahatma Gandhi, Ambedkar or Raja Ram Mohan Roy --- all largely addressed these three issues.
\end{itemize}
\end{KeyConcept}

\begin{KeyConcept}
\begin{itemize}
\item The \textbf{social composition} of the initial reformers had a clear pattern (roughly 95\%): they were \textbf{Western-educated, Hindu, upper-caste, middle-class men}.
\item In short: Western educated + Hindu + upper caste + middle class + men.
\item They shared the three objectives, and were also united by Western education, upper-caste background, and male gender.
\end{itemize}
\end{KeyConcept}

\section{Religious Reform as Social Reform}

\begin{KeyConcept}
\begin{itemize}
\item \textbf{Religious reform implies social reform} because \textbf{religion is a way of life} --- it determines every aspect of social life.
\item In Hinduism, \textbf{caste} defines occupation, marriage, status, loyalty and obligations, whom one can eat with, and even dress and appearance --- i.e. caste defines \textbf{life chances}.
\item Since \textbf{caste is the product of religion}, reforming religion would reform the whole social structure. Hence the early reformers spoke of \textbf{religious reform}, which effectively meant social reform.
\end{itemize}
\end{KeyConcept}

\section{Emancipation of Women: The Status Before Reform}

\begin{KeyConcept}
\begin{itemize}
\item The status of women was \textbf{inferior to men}, and this inferiority had \textbf{religious sanction} (justified in the name of the Vedas, the Quran, etc.). The key facts to note:
\item \textbf{Hindu women could marry only once}, while \textbf{Hindu men could marry multiple times} and keep more than one wife --- polygamy was practised among Hindu men, not only Muslims.
\item \textbf{Prohibition on widow remarriage} applied to women but not to men.
\item \textbf{Muslim women had to observe purdah} (not Muslim men).
\item \textbf{Sati} --- self-immolation on the husband's death (the widow was expected to jump into the fire).
\item \textbf{No right to property} for women.
\item \textbf{No right to divorce/terminate an undesirable marriage} (for both Hindu and Muslim women).
\item \textbf{Social and economic dependence on men} --- women could not go out to work, and their personality was reduced to obeying the husband, cooking, cleaning and caring for the elderly. In the language of \textbf{Ann Oakley}, she was \textbf{alienated from her creative self}.
\item \textbf{Child marriage} --- around age 8--10, largely among upper-caste women (because upper castes were most concerned about purity).
\end{itemize}
\end{KeyConcept}

\begin{ExampleBox}
\textbf{Upper-class vs peasant women}: the \textbf{peasant (largely Shudra) woman} was a \textbf{working woman} --- she could go out, earn, and had more \textbf{agency/mobility} than the upper-caste woman, who was wholly dependent on men. So peasant women had greater agency than upper-caste women.
\end{ExampleBox}

\section{Women's Reform and Its Constitutional Culmination}

\begin{KeyConcept}
\begin{itemize}
\item Social reform attacked these orthodox customs, beginning with the \textbf{Brahmo Samaj}. \textbf{Primary westernization} infused modern practices of equality and humanitarianism across the masses.
\item Key initiators: \textbf{William Bentinck} (abolition of sati) and \textbf{Raja Ram Mohan Roy} (abolition of sati and child marriage, and support for widow remarriage).
\item Gradually the national movement saw large participation of women, culminating in the \textbf{All India Women's Conference (1927)}, founded by \textbf{Margaret Cousins} --- the first all-India body of its kind (like the Indian National Congress, 1885, at the political level). Its objective was to \textbf{promote education and literacy among women at mass level}.
\item Today the body has \textbf{adapted to new problems} --- human trafficking, rape, domestic violence, health and fertility issues --- functioning as an NGO.
\item Other milestones: \textbf{Sarojini Naidu} became the first woman president of the Indian National Congress.
\end{itemize}
\end{KeyConcept}

\begin{GovtPolicy}
\begin{itemize}
\item Reform ideals ultimately culminated in the \textbf{Constitution}: if religion first justified women's inferior status, the Constitution restored their power. Key provisions:
\item \textbf{Articles 14 \& 15} --- equal status for women at par with men (gender equality).
\item \textbf{Directive Principles of State Policy} --- e.g. equal pay for equal work.
\item \textbf{Hindu Succession Act, 1956} --- women became \textbf{co-heirs} to property.
\item \textbf{Hindu Marriage Act, 1955} --- empowered women to dissolve the marriage/seek divorce on certain grounds, and made \textbf{monogamy compulsory} for Hindus.
\item Yet some practices, like \textbf{dowry}, still continue.
\end{itemize}
\end{GovtPolicy}

\section{The Sociological Reading: Private to Public Patriarchy}

\begin{KeyConcept}
\begin{itemize}
\item Using \textbf{Sylvia Walby's} framework: the pre-reform status of women was \textbf{private patriarchy}; today, although women are legally empowered, patriarchy persists as \textbf{public patriarchy}.
\item \textbf{Macro-structures have changed, but the micro reality is almost the same}: legally, norms of mobility and individualism were introduced, but socially, patriarchal norms and customs still prevail over legal rules.
\item In other words, \textbf{universalism was constitutionally imposed from the top, but particularism still prevails at the bottom.}
\end{itemize}
\end{KeyConcept}

\begin{ExampleBox}
\begin{itemize}
\item Example: the \textbf{Hindu Succession Act} made women co-heirs, but women in the majority of cases still do not have \textbf{land ownership} in agriculture. About \textbf{80\% of agricultural labourers are women}, yet hardly \textbf{3--4\% of them own land}.
\item Without land ownership one is a \textbf{peasant}, not a \textbf{farmer} (a farmer owns land); only farmers can avail government schemes and services (better irrigation, soil health card, fertilisers), so landless women are further deprived of opportunities.
\end{itemize}
\end{ExampleBox}

\section{Abolition of Untouchability: The Status Before Reform}

\begin{KeyConcept}
\begin{itemize}
\item Caste defines \textbf{life chances} (occupation, marriage, status, loyalty, food choices). Those \textbf{without caste} within Hinduism --- the \textbf{avarna} (as against the \textbf{savarna} or 'caste Hindus'; the avarna are the 'pancham'/fifth, the 'broken men', i.e. Dalits) --- have their life chances most heavily regulated and are the most vulnerable.
\item The problems faced by the untouchables before reform:
\item In \textbf{South India}, an untouchable walking the same street had to \textbf{announce his presence} so that even his \textbf{shadow} did not fall on a Brahmin (purity--pollution was practised more stringently in the South than the North).
\item In \textbf{Punjab}, the untouchable had to \textbf{shout aloud} to announce his arrival so everyone could hide --- his mere presence was believed to pollute the environment.
\item A separate \textbf{well} was reserved for untouchables; where no separate tank existed, they drank dirty water from ponds/irrigation canals.
\item They were \textbf{prohibited from entering temples}, so they \textbf{manufactured their own gods} (whom Brahmins did not respect).
\item Their children could \textbf{not study in the same school} as upper-caste children --- separate schools were opened.
\item They could \textbf{not join the police or army} or government services, and were \textbf{restricted to menial occupations} (scavenging, tanning hides and skins --- the leather industry, seen as impure because it involves animal skin).
\end{itemize}
\end{KeyConcept}

\section{Reformers Against Untouchability (Phule, Ambedkar, Periyar)}

\begin{KeyConcept}
\begin{itemize}
\item The abolition of untouchability was initiated by \textbf{Jyotiba Phule, Ambedkar and Gandhi}.
\item \textbf{Jyotiba Phule} --- founded the \textbf{Satya Shodak Samaj}; wrote the famous book \textbf{'Gulamgiri'} (slavery) about Dalits, agricultural slaves and bonded labourers. The Samaj \textbf{attacked the Brahmin claim of superiority}: Phule argued the upper castes claim ownership of land while the low castes actually work it --- an \textbf{inherent contradiction in the relations of production} --- so the land belonged to the low castes. He urged the \textbf{Shudras and Ati-Shudras} (the Dalits) to unite against caste discrimination and propagated caste equality.
\item \textbf{Ambedkar} --- the \textbf{temple-entry movement}; founded the \textbf{All India Scheduled Castes Federation}; published the journal \textbf{'Mook Nayak'}; and founded the \textbf{People's Education Society}.
\item \textbf{Gandhi} --- the \textbf{All India Harijan Sangh (1932)}.
\end{itemize}
\end{KeyConcept}

\begin{KeyConcept}
\textbf{Periyar} (E. V. Ramasamy) --- equally important; later linked to the \textbf{Self-Respect Movement} and the \textbf{SNDP (Sree Narayana Dharma Paripalana) movement}. He was an \textbf{outspoken critic of the Hindu scriptures} (like Ambedkar) --- the \textbf{Bhagavad Gita, Ramayana and Manusmriti} --- because he believed they justify the superiority of the Brahmins.
\end{KeyConcept}

\section{Westernization and the Breakdown of the Caste--Occupation Link}

\begin{KeyConcept}
\begin{itemize}
\item Before westernization, \textbf{caste decided occupation} --- an example of \textbf{cumulative inequality}. After westernization, \textbf{meritocracy, egalitarianism and equal opportunity} arrived: your ability/merit (e.g. UPSC marks) decides your occupation --- \textbf{dispersed inequality}. This is \textbf{Andre Beteille's} idea of the \textbf{breakdown of cumulative inequality}.
\item \textbf{Westernization} introduced modern industries, railways and meritocracy, promoting rapid \textbf{urbanization}, increased \textbf{social interaction} and \textbf{intermixing of castes}, and the \textbf{breakdown of the caste--occupation link}.
\item On railways, people of different castes could no longer avoid sitting together or eating train/flight food (you never know who prepared it).
\item With leather becoming a luxury (e.g. \textbf{Louis Vuitton}), even upper-caste merchants and Brahmins entered the leather/hide trade --- \textbf{profit motive} displaced purity--pollution.
\item \textbf{Land reforms} broke the link between caste and land ownership (land ceiling, fragmentation).
\item \textbf{Modern Western education} (the university system) did not discriminate between upper and lower castes, and \textbf{equality before law} took away the judicial functions of the caste panchayat, opening administrative services to all castes.
\end{itemize}
\end{KeyConcept}

\begin{KeyConcept}
\begin{itemize}
\item \textbf{Industrial society rewards instrumental action} (goal-oriented social action) --- the dominant norm became \textbf{profit motive, not purity--pollution} --- responsible for the breakdown of the \textbf{caste--vocation link}.
\item The education system became \textbf{secular and meritocratic}, and \textbf{equality before law} opened the doors of administrative services to all castes.
\end{itemize}
\end{KeyConcept}

Overall, social reform came with westernization --- largely \textbf{heterogenetic changes}, though some were \textbf{orthogenetic}. The reformers (who first contacted westernization --- primary westernization, absorbing rationality, secularism, humanism, equality) built a \textbf{civil society} that spread these ideas to the masses (secondary westernization).

\section{Conclusion: From Particularism to Universalism}

\begin{KeyConcept}
\begin{itemize}
\item Social reform marks the movement from \textbf{traditional society to modern society}:
\item Earlier the norm was \textbf{particularism}; now it is \textbf{universalism}.
\item Earlier \textbf{ascribed/ascriptive features} decided life chances; now \textbf{status is achieved}.
\item Earlier there was \textbf{collective orientation} (collectivism --- the individual, e.g. a Dalit, had to obey the caste structure and Brahmins); now there is \textbf{self-orientation} (every person is equal and can assert rights and hold the government accountable).
\item In Parsons' terms, this is the shift from \textbf{pattern variable A to pattern variable B}, achieved through social reform.
\end{itemize}
\end{KeyConcept}

\begin{KeyConcept}
\textbf{Conclusion}: social reforms introduced the ideas of the \textbf{Renaissance, Enlightenment and the French Revolution} into Indian society. Social reforms \textbf{propelled modernization}, but \textbf{certain traditions have not yet withered away} --- and one can be hopeful that what the social reformers envisioned will be practised and established by \textbf{legal practitioners and state administrators}.
\end{KeyConcept}

\begin{QuickRevision}
\begin{itemize}
\item Social reform (1830s onwards) centred on caste practices, women's status, untouchability.
\item Three objectives: emancipation of women; abolition of untouchability; reform of caste structure.
\item Reformers' social composition: Western-educated, Hindu, upper-caste, middle-class men.
\item Religious reform = social reform (religion/caste define life chances).
\item Women's status: one marriage (vs male polygamy), no widow remarriage, purdah, sati, no property/divorce, child marriage, dependence --- private patriarchy (Walby); reform $\rightarrow$ Brahmo Samaj, Bentinck, Ram Mohan Roy $\rightarrow$ All India Women's Conference 1927 (Margaret Cousins) $\rightarrow$ Articles 14/15, DPSP, Hindu Succession Act 1956, Hindu Marriage Act 1955; but dowry persists, public patriarchy, macro change/micro continuity.
\item Untouchability: avarna/'broken men' most vulnerable; shadow-pollution (South), shout (Punjab), separate wells, no temple entry, separate schools, menial jobs (scavenging, tanning).
\item Reformers: Phule (Satya Shodak Samaj, Gulamgiri, relations of production), Ambedkar (temple entry, SC Federation, Mook Nayak, People's Education Society), Gandhi (Harijan Sangh 1932), Periyar (Self-Respect, SNDP, critique of scriptures).
\item Westernization: cumulative $\rightarrow$ dispersed inequality (Beteille); railways/meritocracy/land reforms/secular education broke caste--occupation link; profit motive > purity--pollution.
\item Conclusion: particularism $\rightarrow$ universalism; ascribed $\rightarrow$ achieved; collectivism $\rightarrow$ self-orientation; Renaissance/Enlightenment/French Revolution ideas; modernization with persisting traditions.
\end{itemize}
\end{QuickRevision}

